\chapter{外部数据及其约束范围}
\label{ch:data}

公开数据只约束它测到的量。一条等温线给出的是平衡含水与活度、温度的关系。它不给出冷点轨迹，不给出探针的信息增益，也不给出蒸烤过程中的非平衡含水。本章记录已经读入的拟合表，以及尚缺、因而不能写入主结果的对象。

\section{已读入的等温线}

\paragraph{要判定的量}
清洁拟合能收紧哪一个先验，以及这个先验的温度上限停在哪里。

\paragraph{继承的文献层与不继承的层}
GAB 的三参数形式继承多层吸附的结构 \cite{gab1980}。参数数值来自已经读入的拟合表 \cite{sorption2026}，许可为 CC BY 4.0，记录号 \texttt{10.5281/zenodo.21671925}。不继承把拟合值写成新的实验测定。不继承把该表没有的肉类曲线补成“文献中不存在”。未读到不等于不存在。

读入的是工作表 Original GAB on ln(Mc)。全部曲线记录 200 条。模型码为 GAB2 且诊断码为 0 的清洁拟合 140 条。诊断码不为 0 的 60 条不进入先验：其中有拟合失败，也有 \(K\) 被限制到 1 之后退化为 BET 的行。另一张工作表是重参数化拟合，本次先验不用它，避免同一曲线以两套参数进入区间。

在 140 条上，按排序样本线性插值，\(X_m\) 的十分位、中位、九十分位为 \(0.0366\)、\(0.0733\)、\(0.1453\)；\(C\) 为 \(3.98\)、\(11.20\)、\(45.24\)；\(K\) 为 \(0.607\)、\(0.773\)、\(0.959\)。这些分位描述该清洁子集，不描述全部食品。

类别以谷物、果蔬和粉体为主。按记录条数，清洁子集里较多的是稻米 22 条、苋菜 21 条、啤酒糟 19 条、番茄 14 条、乳粉 9 条。按名称检索肉类、牛肉、猪肉、鸡肉和鱼，条数为 0。

土豆有 6 条清洁拟合，温度只到 \(60^\circ\mathrm{C}\)。把式 \eqref{eq:gab} 代进吸附 \(30^\circ\mathrm{C}\) 的参数 \(X_m=0.0567\)、\(C=14.34\)、\(K=0.912\)，活度 \(0.2\)、\(0.4\)、\(0.6\)、\(0.8\) 对应的干基含水约为 \(0.053\)、\(0.080\)、\(0.118\)、\(0.204\)。同一活度下，吸附 \(60^\circ\mathrm{C}\) 的参数 \(X_m=0.0359\)、\(C=6.38\)、\(K=0.762\) 给出 \(0.023\)、\(0.038\)、\(0.056\)、\(0.084\)。温度升高时，同一活度能容纳的水减少。方向与吸附热随温度下降一致。\(60^\circ\mathrm{C}\) 仍低于烤制表面，数值不能外推到表面褐变温度。

\begin{table}[ht]
\centering
\caption{这份等温线能够约束和不能约束的对象}
\begin{tabularx}{\textwidth}{lX}
\toprule
能 & 不能 \\
\midrule
淀粉类 \(a_w(M,T)\) 在测定温度内的先验区间 & 肉类、脂肪组织的等温线 \\
GAB 三参数的数量级与相关结构 & 蒸烤过程中的动态非平衡含水 \\
温度升高时 \(X_m\) 下降这一方向 & \(100^\circ\mathrm{C}\) 以上的外推数值 \\
灵敏度扫描的参数上下界 & 冷点、探针信息增益、控制效果 \\
\bottomrule
\end{tabularx}
\end{table}

使用时保留记录号、类别、温度、吸附或解吸、诊断码和参数。诊断码不为 0 的行不进入先验。

\subsection{同一组土豆参数的正算与反算}

\eqref{eq:gab} 在固定温度上对活度单调。极点在 \(a_w=1/K\)。\(30^\circ\mathrm{C}\) 的 \(K=0.912\)，极点约 \(1.096\)；\(60^\circ\mathrm{C}\) 的 \(K=0.762\)，极点约 \(1.312\)。干基含水在极点之前有上界。

\begin{table}[ht]
\centering
\caption{土豆吸附支的正算。参数是该温度上的清洁拟合，不是核的常数。}
\begin{tabular}{lrr}
\toprule
\(a_w\) & \(30^\circ\mathrm{C}\) 干基 & \(60^\circ\mathrm{C}\) 干基 \\
\midrule
\(0.20\) & \(0.0528\) & \(0.0226\) \\
\(0.40\) & \(0.0796\) & \(0.0380\) \\
\(0.60\) & \(0.1184\) & \(0.0558\) \\
\(0.80\) & \(0.2044\) & \(0.0836\) \\
\(0.90\) & \(0.3117\) & \(0.1066\) \\
\(0.95\) & \(0.4199\) & \(0.1227\) \\
\(0.99\) & \(0.5795\) & \(0.1391\) \\
\bottomrule
\end{tabular}
\end{table}

反算把干基含水代回活度。\(30^\circ\mathrm{C}\)：干基 \(0.05\)、\(0.10\)、\(0.20\)、\(0.40\) 对应活度 \(0.180\)、\(0.519\)、\(0.794\)、\(0.943\)。即使活度取到 \(0.999\)，干基仍是 \(0.633\)，到不了 \(0.72\)。\(60^\circ\mathrm{C}\)：干基 \(0.05\) 对应 \(0.542\)，\(0.10\) 对应 \(0.875\)；干基 \(0.20\) 及以上没有根。活度 \(0.999\) 时干基只有 \(0.143\)。

核里的含水 \(0.72\) 是湿基式的分数，上限裁到 \(0.95\)，不是这张表的干基。把核在含水 \(0.10\) 上的土豆活度 \(0.86\) 与 \(30^\circ\mathrm{C}\) 反算的 \(0.52\) 并列，得到的是两套定义的差，不是拟合残差。S1 的三档初含水都在 \(0.64\) 以上，落在这张干基表的值域之外。

\paragraph{本节不能推出的结论}
土豆六条不能写成肉糜的活度测定。核里的活度映射是含水的单调饱和函数，不是 \eqref{eq:gab} 的反函数，而且不随温度变化。因此即使淀粉先验已经读入，本次轨迹也没有使用这六条的 \(X_m(T)\)。干基 \(0.72\) 没有根，也不能写成核的含水设置是错的：两边不是同一个含水定义。

\section{尚缺的测量与仿真如何补位}

\paragraph{要判定的量}
缺测对象在测量到达之前用什么临时约束，到达之后替换成什么。临时约束不能冒充测定。

下列对象在本次使用的公开拟合表里没有对应文件。肌肉与脂肪的导热、烤制中心温度历程、腔体换热曲线和微生物热失活的同步进度，都不在该表中。

\begin{table}[ht]
\centering
\caption{缺测对象的临时约束}
\begin{tabularx}{\textwidth}{lXl}
\toprule
对象 & 在测量到达前 & 到达后替换成 \\
\midrule
肌肉与脂肪导热 & 文献区间；核内取 \(k=0.68\)，脂肪控制体乘 \(0.45\) & 该文件的温度—含水表 \\
内部温度与质量历程 & 只在同一核的策略之间对照 & 外部残差核对 \\
腔体换热与混合 & \(h=12+55u_{\mathrm{fan}}\)，标记为设定边界 & 实测 \(h(v)\) \\
反应进度 & 动力学文献的数量级；核内 \(F\) 的除数为 2 & 同步测量的进度 \\
微生物绝对初始量 & 不进入状态 & 仍不进入状态 \\
\bottomrule
\end{tabularx}
\end{table}

仿真补的是因子覆盖，不是测量。S1 的因子是 3 种食材、3 档厚度、2 个初温、3 档初含水、3 档非均质、30 个种子，共 \(3\times 3\times 2\times 3\times 3\times 30=4860\) 个组合。四种策略各跑一遍，共 19440 行。早期提纲里的 1620 没有把非均质的三档乘进去。日志以 4860 为准，不把行数改回 1620，也不把 19440 写成 19440 次烹饪实验。

参数文件的来源字段只允许三个值：\texttt{doi:10.5281/zenodo.21671925}、\texttt{literature-interval}、\texttt{simulation}。第三个值不出现在摘要的验证句里。本次战役每一行的来源字段都是 \texttt{simulation}，扩充标记都是假。因此摘要里不出现“由公开等温线验证了冷点”这类句子。

后续测量文件一旦取得，先按下一节的入口核对，再决定是否缩小该参数的抽样区间。不重跑全部格子，除非核对残差改变了按厚度分开的策略排序。

\paragraph{本节不能推出的结论}
临时约束写在表里，不能写成该导热已经由 \cite{choi1986} 的某一行测定。各向异性不在组分加和式内。

\section{数据进入模型的唯一入口}

\paragraph{要判定的量}
外部测量接触在线后验的位置。若有两个以上入口，同一参数会在先验和似然里被用两次。

入口只有两处。

\begin{enumerate}
\item \textbf{先验区间。} 清洁 GAB 参数写成淀粉类别的上下界。在线更新不得跑出该界，除非本次质量与湿度残差证明界本身与样品矛盾，并在日志中标记越界。肉类没有清洁行，肉类活度保持为开放区间，不套用土豆的 \(X_m\)。
\item \textbf{外部核对。} 若以后的公开温度历程与仿真边界足够接近，比较的是 \(T(t)\) 与质量损失的残差。残差过大时收缩几何或边界的适用范围，不改命题 \ref{prop:steady}--\ref{prop:shield} 的陈述。那些命题依赖的是方程结构，不是某一条等温线的拟合值。
\end{enumerate}

核里的活度函数把原型收成两个常数 \((s_p,k_p)\)：肉糜 \((0.10,0.70)\)、土豆 \((0.055,0.80)\)、面团 \((0.08,0.75)\)。在 S1 的初含水 \(\{0.64,0.72,0.80\}\) 上，三个原型的活度都饱和到 \(0.99\) 以上。差别要等表面含水降到大约 \(0.3\) 以下才显现：含水 \(0.20\) 时肉糜活度约 \(0.88\)，土豆约 \(0.98\)，面团约 \(0.93\)。因此食材交叉表在停机误差上几乎重合，与“高含水段活度饱和”一致。这不是蛋白反应和淀粉反应已经被检验。

演示下载的参数文件要带来源字段。无来源的肉类等温线不得标成来自这份数据。

下面列出全部清洁拟合，供先验区间核对。

\subsection{清洁拟合的记录}
下列数字读自工作表 Original GAB on ln(Mc)，模型码为 GAB2，诊断码为 0。该表全部曲线记录 200 条，清洁拟合 140 条。三个参数是单层容量 $X_m$、能量常数 $C$ 和多层校正 $K$。
清洁子集上 $X_m$ 的十分位、中位、九十分位为 0.0366、0.0733、0.1453；$C$ 为 3.98、11.20、45.24；$K$ 为 0.607、0.773、0.959。分位按排序样本的线性插值计算，只描述这 140 条，不描述全部食品。
\begin{table}[ht]
\centering
\caption{土豆清洁拟合，温度只到 $60^\circ\mathrm{C}$}
\begin{tabular}{llrrrrrr}
\toprule
记录 & 方向 & $T/^\circ\mathrm{C}$ & $X_m$ & $C$ & $K$ & $n$ & $R^2$ \\
\midrule
c074 & ADS & 30 & 0.0567 & 14.34 & 0.912 & 7 & 0.9963 \\
c075 & ADS & 45 & 0.0457 & 13.71 & 0.885 & 7 & 0.9970 \\
c076 & ADS & 60 & 0.0359 & 6.38 & 0.762 & 7 & 0.9990 \\
c077 & DES & 30 & 0.0763 & 9.21 & 0.861 & 7 & 0.9932 \\
c078 & DES & 45 & 0.0549 & 8.55 & 0.860 & 7 & 0.9975 \\
c079 & DES & 60 & 0.0522 & 3.75 & 0.666 & 7 & 0.9984 \\
\bottomrule
\end{tabular}
\end{table}
肉类名称在清洁子集中的条数为 0。谷物、果蔬和粉体占全部 140 条。这些行约束淀粉类活度在测定温度内的数量级。它们不约束冷点、探针信息增益和烤制表面。
\begin{longtable}{llrrrrrr}
\caption{诊断码为 0 的 GAB2 拟合}\\
\toprule
记录 & 类别 & $T/^\circ\mathrm{C}$ & $X_m$ & $C$ & $K$ & $n$ & $R^2$ \\
\midrule
\endfirsthead
\toprule
记录 & 类别 & $T/^\circ\mathrm{C}$ & $X_m$ & $C$ & $K$ & $n$ & $R^2$ \\
\midrule
\endhead
c012 & milkPowder & 5 & 0.0720 & 4.21 & 0.859 & 5 & 0.988 \\
c013 & milkPowder & 15 & 0.0502 & 4.92 & 0.904 & 6 & 0.994 \\
c014 & milkPowder & 15 & 0.0510 & 6.76 & 0.918 & 5 & 0.998 \\
c015 & milkPowder & 20 & 0.0409 & 6.66 & 0.957 & 6 & 0.994 \\
c017 & milkPowder & 25 & 0.0375 & 7.19 & 0.944 & 6 & 0.999 \\
c019 & milkPowder & 30 & 0.0353 & 7.14 & 0.936 & 6 & 1.000 \\
c020 & milkPowder & 30 & 0.0404 & 7.76 & 0.926 & 5 & 0.999 \\
c021 & milkPowder & 35 & 0.0293 & 8.58 & 0.962 & 6 & 0.996 \\
c022 & milkPowder & 35 & 0.0384 & 7.95 & 0.911 & 5 & 0.997 \\
c028 & tomato & 20 & 0.2037 & 27.48 & 0.834 & 7 & 0.998 \\
c029 & tomato & 20 & 0.1940 & 29.04 & 0.853 & 7 & 0.997 \\
c030 & tomato & 20 & 0.2074 & 25.97 & 0.826 & 7 & 0.999 \\
c031 & tomato & 30 & 0.1621 & 22.65 & 0.897 & 7 & 0.985 \\
c032 & tomato & 30 & 0.1601 & 21.48 & 0.905 & 7 & 0.986 \\
c033 & tomato & 30 & 0.1586 & 24.67 & 0.900 & 7 & 0.986 \\
c034 & tomato & 40 & 0.1610 & 10.81 & 0.896 & 7 & 0.991 \\
c035 & tomato & 40 & 0.1648 & 9.29 & 0.888 & 7 & 0.987 \\
c036 & tomato & 40 & 0.1593 & 12.74 & 0.902 & 7 & 0.988 \\
c037 & tomato & 50 & 0.2002 & 4.43 & 0.840 & 7 & 0.999 \\
c038 & tomato & 50 & 0.2090 & 4.01 & 0.832 & 7 & 0.993 \\
c039 & tomato & 50 & 0.2247 & 3.53 & 0.809 & 7 & 0.999 \\
c040 & tomato & 60 & 0.1097 & 10.22 & 0.972 & 7 & 0.959 \\
c041 & tomato & 60 & 0.1165 & 9.36 & 0.960 & 7 & 0.983 \\
c047 & beetroot & 50 & 0.1169 & 1.65 & 0.926 & 9 & 0.990 \\
c048 & beetroot & 50 & 0.1274 & 1.32 & 0.937 & 9 & 0.986 \\
c074 & potato & 30 & 0.0567 & 14.34 & 0.912 & 7 & 0.996 \\
c075 & potato & 45 & 0.0457 & 13.71 & 0.885 & 7 & 0.997 \\
c076 & potato & 60 & 0.0359 & 6.38 & 0.762 & 7 & 0.999 \\
c077 & potato & 30 & 0.0763 & 9.21 & 0.861 & 7 & 0.993 \\
c078 & potato & 45 & 0.0549 & 8.55 & 0.860 & 7 & 0.998 \\
c079 & potato & 60 & 0.0522 & 3.75 & 0.666 & 7 & 0.998 \\
c081 & rice & 5 & 0.0697 & 5.85 & 0.787 & 9 & 0.995 \\
c083 & rice & 15 & 0.0778 & 2.97 & 0.772 & 9 & 0.994 \\
c084 & rice & 20 & 0.0679 & 2.65 & 0.809 & 9 & 0.995 \\
c085 & rice & 25 & 0.0645 & 5.86 & 0.793 & 9 & 0.996 \\
c086 & rice & 30 & 0.0751 & 3.11 & 0.754 & 9 & 0.993 \\
c088 & mushroom & 30 & 0.1722 & 4.72 & 0.713 & 5 & 0.999 \\
c092 & mushroom & 50 & 0.1391 & 4.70 & 0.710 & 5 & 0.993 \\
c093 & mushroom & 50 & 0.1108 & 3.11 & 0.676 & 5 & 0.991 \\
c094 & mushroom & 60 & 0.1676 & 3.31 & 0.604 & 5 & 0.995 \\
c095 & mushroom & 60 & 0.0677 & 4.94 & 0.765 & 5 & 0.989 \\
c097 & mushroom & 70 & 0.0626 & 3.45 & 0.740 & 5 & 0.983 \\
c098 & flour & 10 & 0.0823 & 15.13 & 0.722 & 9 & 0.999 \\
c099 & flour & 20 & 0.0801 & 13.48 & 0.734 & 9 & 0.999 \\
c100 & flour & 30 & 0.0781 & 12.01 & 0.740 & 9 & 0.999 \\
c101 & millet & 20 & 0.0743 & 12.82 & 0.796 & 6 & 0.994 \\
c102 & millet & 30 & 0.0673 & 10.55 & 0.806 & 6 & 0.995 \\
c103 & millet & 40 & 0.0650 & 9.92 & 0.794 & 6 & 1.000 \\
c104 & millet & 50 & 0.0545 & 11.93 & 0.778 & 6 & 0.992 \\
c106 & millet & 30 & 0.0758 & 26.98 & 0.791 & 6 & 0.985 \\
c108 & millet & 50 & 0.0571 & 35.91 & 0.837 & 6 & 0.995 \\
c109 & milet & 20 & 0.0801 & 9.04 & 0.774 & 6 & 0.982 \\
c110 & milet & 30 & 0.0690 & 7.72 & 0.774 & 6 & 0.989 \\
c111 & milet & 40 & 0.0667 & 7.98 & 0.772 & 6 & 0.994 \\
c112 & milet & 50 & 0.0448 & 11.22 & 0.820 & 6 & 0.998 \\
c113 & milet & 20 & 0.0984 & 10.82 & 0.706 & 6 & 0.996 \\
c114 & milet & 30 & 0.0946 & 7.37 & 0.693 & 6 & 0.971 \\
c115 & milet & 40 & 0.0788 & 8.23 & 0.734 & 6 & 0.977 \\
c116 & milet & 50 & 0.0824 & 3.94 & 0.701 & 6 & 0.980 \\
c117 & rice & 10 & 0.0989 & 24.36 & 0.685 & 4 & 1.000 \\
c118 & rice & 20 & 0.1035 & 17.06 & 0.632 & 4 & 0.984 \\
c119 & rice & 30 & 0.1108 & 12.30 & 0.588 & 4 & 1.000 \\
c120 & rice & 40 & 0.0940 & 10.19 & 0.674 & 4 & 1.000 \\
c122 & rice & 10 & 0.0855 & 25.16 & 0.737 & 4 & 1.000 \\
c123 & rice & 20 & 0.0878 & 16.93 & 0.702 & 4 & 0.985 \\
c124 & rice & 30 & 0.0899 & 12.81 & 0.679 & 4 & 1.000 \\
c125 & rice & 40 & 0.0830 & 9.26 & 0.717 & 4 & 0.999 \\
c127 & oat & 5 & 0.0291 & 64.42 & 0.943 & 8 & 0.975 \\
c128 & oat & 20 & 0.0311 & 23.23 & 0.925 & 8 & 0.994 \\
c129 & oat & 40 & 0.0319 & 11.24 & 0.901 & 8 & 0.992 \\
c130 & oat & 60 & 0.0277 & 10.74 & 0.907 & 8 & 0.997 \\
c131 & amaranth & 25 & 0.0535 & 6.75 & 0.768 & 8 & 0.998 \\
c132 & amaranth & 25 & 0.0728 & 17.38 & 0.657 & 8 & 0.997 \\
c133 & amaranth & 30 & 0.0555 & 16.54 & 0.790 & 8 & 0.998 \\
c134 & amaranth & 30 & 0.0783 & 21.47 & 0.688 & 8 & 0.998 \\
c135 & amaranth & 35 & 0.0548 & 10.11 & 0.752 & 8 & 0.992 \\
c136 & amaranth & 35 & 0.0675 & 18.76 & 0.686 & 8 & 0.993 \\
c137 & amaranth & 35 & 0.0695 & 10.13 & 0.768 & 10 & 0.987 \\
c138 & amaranth & 40 & 0.0776 & 8.31 & 0.545 & 7 & 0.997 \\
c139 & amaranth & 40 & 0.0782 & 14.46 & 0.566 & 7 & 0.996 \\
c140 & amaranth & 40 & 0.0749 & 10.43 & 0.763 & 4 & 0.995 \\
c141 & amaranth & 45 & 0.0738 & 7.64 & 0.587 & 5 & 0.999 \\
c142 & amaranth & 45 & 0.0705 & 13.58 & 0.668 & 7 & 0.998 \\
c143 & amaranth & 45 & 0.0604 & 10.17 & 0.791 & 9 & 0.987 \\
c144 & amaranth & 50 & 0.0576 & 10.44 & 0.715 & 7 & 0.991 \\
c145 & amaranth & 50 & 0.0696 & 14.25 & 0.623 & 7 & 0.999 \\
c146 & amaranth & 50 & 0.0791 & 6.10 & 0.720 & 4 & 0.999 \\
c147 & amaranth & 55 & 0.0498 & 7.16 & 0.783 & 7 & 0.996 \\
c148 & amaranth & 55 & 0.0601 & 8.75 & 0.698 & 7 & 0.996 \\
c149 & amaranth & 65 & 0.0593 & 7.31 & 0.724 & 7 & 0.992 \\
c150 & amaranth & 70 & 0.0820 & 4.71 & 0.678 & 4 & 0.996 \\
c151 & amaranth & 90 & 0.0893 & 2.77 & 0.651 & 4 & 1.000 \\
c152 & rice & 40 & 0.1352 & 5.79 & 0.476 & 8 & 1.000 \\
c153 & rice & 50 & 0.0967 & 5.82 & 0.608 & 8 & 0.999 \\
c154 & rice & 60 & 0.1166 & 3.98 & 0.531 & 8 & 0.996 \\
c155 & rice & 70 & 0.1306 & 2.94 & 0.504 & 9 & 0.995 \\
c156 & rice & 80 & 0.1438 & 2.25 & 0.504 & 9 & 0.996 \\
c157 & corn & 20 & 0.0743 & 32.55 & 0.700 & 10 & 0.994 \\
c158 & corn & 25 & 0.0727 & 28.61 & 0.696 & 10 & 0.977 \\
c159 & corn & 40 & 0.0673 & 15.66 & 0.725 & 10 & 0.997 \\
c160 & corn & 20 & 0.0996 & 26.57 & 0.591 & 10 & 0.981 \\
c161 & corn & 25 & 0.0963 & 23.82 & 0.607 & 10 & 0.988 \\
c162 & corn & 40 & 0.0943 & 17.15 & 0.617 & 10 & 0.991 \\
c163 & rice & 40 & 0.0830 & 12.58 & 0.677 & 8 & 0.989 \\
c164 & rice & 50 & 0.0872 & 8.04 & 0.589 & 6 & 0.998 \\
c165 & rice & 60 & 0.0567 & 9.52 & 0.715 & 10 & 0.998 \\
c166 & rice & 70 & 0.0461 & 6.39 & 0.756 & 8 & 0.978 \\
c167 & wheat & 45 & 0.1001 & 30.98 & 0.597 & 6 & 0.992 \\
c168 & wheat & 60 & 0.0492 & 11.77 & 0.917 & 12 & 0.996 \\
c169 & sorgum & 20 & 0.0773 & 37.22 & 0.670 & 10 & 0.984 \\
c170 & sorgum & 25 & 0.0746 & 32.14 & 0.687 & 10 & 0.989 \\
c171 & sorgum & 40 & 0.0721 & 14.60 & 0.697 & 10 & 0.992 \\
c172 & sorgum & 20 & 0.0982 & 28.06 & 0.598 & 10 & 0.980 \\
c173 & sorgum & 25 & 0.0961 & 27.52 & 0.608 & 10 & 0.979 \\
c174 & sorgum & 40 & 0.0936 & 19.18 & 0.614 & 10 & 0.991 \\
c175 & quinoa & 20 & 0.0893 & 19.60 & 0.667 & 7 & 0.932 \\
c176 & quinoa & 30 & 0.0884 & 11.10 & 0.660 & 10 & 0.998 \\
c177 & quinoa & 40 & 0.0552 & 8.42 & 0.833 & 10 & 0.962 \\
c178 & quinoa & 20 & 0.1073 & 17.53 & 0.558 & 7 & 0.915 \\
c179 & quinoa & 30 & 0.0890 & 11.18 & 0.660 & 10 & 0.999 \\
c180 & quinoa & 40 & 0.0501 & 12.12 & 0.871 & 9 & 0.962 \\
c181 & brewerSpentGrain & 5 & 0.0589 & 39.69 & 0.937 & 14 & 1.000 \\
c182 & brewerSpentGrain & 5 & 0.0566 & 25.73 & 0.937 & 14 & 0.998 \\
c183 & brewerSpentGrain & 10 & 0.0481 & 56.81 & 0.952 & 14 & 0.998 \\
c184 & brewerSpentGrain & 10 & 0.0514 & 47.34 & 0.938 & 14 & 0.997 \\
c185 & brewerSpentGrain & 20 & 0.0573 & 58.60 & 0.946 & 14 & 0.999 \\
c186 & brewerSpentGrain & 20 & 0.0455 & 35.93 & 0.959 & 14 & 0.998 \\
c187 & brewerSpentGrain & 30 & 0.0426 & 49.66 & 0.959 & 14 & 0.997 \\
c188 & brewerSpentGrain & 30 & 0.0404 & 46.69 & 0.956 & 14 & 0.998 \\
c189 & brewerSpentGrain & 40 & 0.0412 & 46.40 & 0.962 & 13 & 0.998 \\
c190 & brewerSpentGrain & 40 & 0.0366 & 61.10 & 0.965 & 13 & 0.998 \\
c191 & brewerSpentGrain & 50 & 0.0384 & 54.06 & 0.965 & 13 & 0.998 \\
c192 & brewerSpentGrain & 50 & 0.0364 & 47.63 & 0.964 & 13 & 0.997 \\
c193 & brewerSpentGrain & 60 & 0.0376 & 50.79 & 0.970 & 9 & 0.999 \\
c194 & brewerSpentGrain & 60 & 0.0346 & 56.01 & 0.976 & 9 & 0.999 \\
c195 & brewerSpentGrain & 70 & 0.0354 & 47.36 & 0.981 & 9 & 0.999 \\
c196 & brewerSpentGrain & 70 & 0.0326 & 40.52 & 0.974 & 9 & 0.996 \\
c197 & brewerSpentGrain & 80 & 0.0338 & 43.83 & 0.971 & 8 & 0.999 \\
c198 & brewerSpentGrain & 80 & 0.0336 & 45.11 & 0.952 & 8 & 0.999 \\
c200 & brewerSpentGrain & 90 & 0.0312 & 62.84 & 0.973 & 6 & 0.995 \\
\bottomrule
\end{longtable}
表中类别名称保持数据文件中的拼写。诊断码不为 0 的 60 条不进入先验：其中包含拟合失败，以及 $K$ 被限制到 1 之后退化为 BET 的行。


\paragraph{本节不能推出的结论}
140 条全部列出，不能推出其中任意一条已经进入本次积分器。积分器的活度不随温度变化，也不读取记录号。
